%%% ============================================================
%%% The Discrete Fejer Quotient on Squarefree Moduli:
%%%   Spectral Characterization, Layered Divisors, and the
%%%   Asymptotic Corridor Average
%%%
%%% B. R. Sanders, M. Gish (2026)
%%% Target venue: Integers -- Electronic Journal of Combinatorial
%%%               Number Theory
%%% Backup venue: AMM Notes, Math. Magazine, Math. Intelligencer
%%% MSC: 11A41, 11N05, 11A51, 11Y05, 42A16
%%% DOI of bundle: 10.5281/zenodo.18852047
%%%
%%% Merge history (2026-05-13): this manuscript consolidates the
%%% previously separate J24 ("The First-G Event and a Discrete
%%% Sinc^2 Identity") and J41 ("Full-Period Cancellation of R(k,f)
%%% and the spf-Localization for Squarefree Moduli") into a single
%%% Integers submission. Both prior manuscripts addressed the
%%% zero structure of the discrete Fejer quotient R(k,f) on
%%% squarefree moduli from complementary vantage points (smallest
%%% spectral zero vs full-period cancellation; prime-product vs
%%% divisor-product). The circular citation pattern and the
%%% per-quarter cap concern made consolidation the right move.
%%% The J41 manuscript is preserved per never-delete at
%%%   05_papers/number_theory/J41/manuscript/manuscript.tex
%%% with a clear pointer to this consolidated version.
%%%
%%% Author lane: Sanders + Gish only.
%%% ============================================================

\documentclass[11pt,reqno]{amsart}

\usepackage{amsmath,amssymb,amsthm,mathtools}
\usepackage[margin=1in]{geometry}
\usepackage[colorlinks=true,linkcolor=blue,citecolor=blue,urlcolor=blue]{hyperref}
\usepackage{microtype}
\usepackage{enumitem}
\usepackage{booktabs}

%% -------- theorem environments --------
\theoremstyle{plain}
\newtheorem{theorem}{Theorem}[section]
\newtheorem{lemma}[theorem]{Lemma}
\newtheorem{proposition}[theorem]{Proposition}
\newtheorem{corollary}[theorem]{Corollary}

\theoremstyle{definition}
\newtheorem{definition}[theorem]{Definition}

\theoremstyle{remark}
\newtheorem{remark}[theorem]{Remark}

%% -------- macros --------
\DeclareMathOperator{\spf}{spf}
\DeclareMathOperator{\sinc}{sinc}
\DeclareMathOperator{\rad}{rad}
\DeclareMathOperator{\Si}{Si}
\newcommand{\Z}{\mathbb{Z}}
\newcommand{\N}{\mathbb{N}}
\newcommand{\R}{\mathbb{R}}
\newcommand{\C}{\mathbb{C}}
\newcommand{\Q}{\mathbb{Q}}
\newcommand{\Prm}{\mathbb{P}}

%% -------- title matter --------
\title[The discrete Fej\'{e}r quotient on squarefree moduli]%
{The Discrete Fej\'{e}r Quotient on Squarefree Moduli:\\
 Spectral Characterization, Layered Divisors, and the\\
 Asymptotic Corridor Average}

\author{B.~R.~Sanders}
\address{7Site LLC, Hot Springs, Arkansas, USA}
\email{brayden@7site.co}

\author{M.~Gish}
\address{Independent Researcher, Hot Springs, Arkansas, USA}
\email{monica.gish1992@gmail.com}

\date{\today}

\keywords{discrete Fej\'{e}r quotient, smallest prime factor, coprimality
partition, full-period cancellation, squarefree moduli, layered divisor
lattice, Boolean divisor lattice, sinc function, sine integral, spectral
product, obstruction-zero correspondence}

\subjclass[2020]{11A41, 11N05, 11A51, 11Y05, 42A16}

\begin{document}

\begin{abstract}
For an integer $f \ge 2$ and an integer $k \ge 1$, the
\emph{discrete Fej\'{e}r quotient}
$R(k, f) = \sin^{2}(\pi k/f) / (k^{2} \sin^{2}(\pi/f))$ is the
squared modulus of the $k$-element average of the geometric progression
$e^{2\pi i j/f}$ ($j = 1, \dots, k$). We organize the zero structure of
$R$ on squarefree moduli via three complementary results.
\emph{(i)~Full-period cancellation} (Theorem~\ref{thm:fullperiod}):
$R(k, f) = 0$ if and only if $f \mid k$, uniformly in $f \ge 2$.
\emph{(ii)~Spectral characterization of the obstruction sequence}
(Theorem~\ref{thm:obstruction-zero}): for every integer $b > 1$ with
distinct prime factors $p_{1}, \dots, p_{r}$, the spectral product
$f_{b}(k) := \prod_{j=1}^{r} R(k, p_{j})$ vanishes at $k \in \N$
exactly when $\gcd(k, b) > 1$. The synchronization at the smallest
spectral zero (Theorem~\ref{thm:sync}) localizes the First-G event of
the coprimality partition at $k = \spf(b)$, and the asymptotic zero
density of $f_{b}$ is the Euler product
$1 - \varphi(\rad(b))/\rad(b)$ (Corollary~\ref{cor:asymp-density}).
\emph{(iii)~Layered-divisor structure} (Theorem~\ref{thm:layered}):
for squarefree $b = p_{1} \cdots p_{r}$ and the $j$-th primorial
divisor $b_{j} = p_{1} \cdots p_{j}$, exactly $2^{j} - 1$ non-trivial
divisors $d \mid b$ satisfy $R(b_{j}, d) = 0$. The number-theoretic
content is bridged to the harmonic side: the discrete Fej\'{e}r
quotient converges to the continuum
$\sinc^{2}(k/f)$ as $f \to \infty$ with $k/f$ fixed
(Theorem~\ref{thm:limit}), and the corridor average converges to
the sine-integral constant
$\frac{1}{f-1}\sum_{k=1}^{f-1} R(k, f) \to \Si(2\pi)/\pi
\approx 0.4514$ (Theorem~\ref{thm:asymptotic}). All claims are
verified at exact integer arithmetic or machine precision: 22{,}367
$(b, k)$ pairs for the First-G localization, 4{,}225 $(p, k)$ pairs
for the basic biconditional, 900 cells for the obstruction-zero
correspondence, the Euler-product density matched to within
$7 \times 10^{-6}$ for $b$ up to 2310, and Riemann-sum convergence
to $\Si(2\pi)/\pi$ within $5 \times 10^{-5}$ at $f = 1000$.
\end{abstract}

\maketitle

%% -------- section 0: lens and substrate --------
\section*{Lens and substrate}\label{sec:lens}

This paper works on $\Z$ (no specialized substrate) with the standard
objects: smallest prime factor $\spf(b)$, the radical $\rad(b)$, the
discrete Fej\'{e}r quotient $R(k, f)$, and the continuum
$\sinc^{2}$ function. The objects are not derived from a framework
reading; they are the elementary discrete-Fourier and divisibility
structures of the integers. The substantive contribution is the joint
treatment of three complementary structural results on the zero set of
$R(\cdot, \cdot)$ on squarefree moduli: a uniform-in-$f$ full-period
cancellation theorem, a spectral characterization of the obstruction
sequence by means of a product of Fej\'{e}r quotients, and a layered
closure on the divisor lattice of squarefree $b$. We choose the
harmonic-content presentation (smallest prime factor $\times$ discrete
Fej\'{e}r quotient on $\Z$) over a $\Z/N\Z$-algebraic framing because
the obstruction-zero correspondence is a statement about integers and
their prime divisibility, with the spectral product $f_{b}$ supplying
a single naturally-defined Fourier indicator for $\gcd(k, b) > 1$.
The natural setting is $\Z$, not $\Z/N\Z$ with operator-labeled tables;
no magma structure is required for the present results, and the
divisibility lattice of $\rad(b)$ enters directly through the zero loci
of the factors $R(\cdot, p_{j})$. Companion papers in our wider
program~\cite{SandersGish2026Sigma,SandersGish2026FourCore} work over
$\Z/N\Z$ with operator-labeled tables; the present paper does not
require any of that machinery and stands alone.

%% -------- section 1: introduction --------
\section{Introduction}\label{sec:intro}

The discrete Fej\'{e}r quotient at integer parameters $k \ge 1$ and
$f \ge 2$,
\begin{equation}\label{eq:Rdef}
  R(k, f) \;=\; \frac{\sin^{2}(\pi k/f)}{k^{2}\,\sin^{2}(\pi/f)},
\end{equation}
is the standard specialization of the Fej\'{e}r
kernel~\cite{Fejer1900,Zygmund2002} to rational arguments and arises
as the squared modulus of the $k$-element average of the geometric
progression $e^{2\pi i j/f}$ ($j = 1, \dots, k$):
\begin{equation}\label{eq:Rgeom}
  R(k, f)
  \;=\;
  \left|\,\frac{1}{k}\sum_{j=1}^{k} e^{2\pi i j/f}\,\right|^{2}.
\end{equation}
The closed form~\eqref{eq:Rdef} is elementary, but the zero structure
of $R$ has direct content for the divisibility lattice on $\Z$.
This paper records seven theorems organizing that structure on
squarefree moduli.

\paragraph{The seven theorems.} The closed form for $R$
(Theorem~\ref{thm:closedform}) gives the underlying identity. The
full-period cancellation $R(k, f) = 0 \Leftrightarrow f \mid k$
(Theorem~\ref{thm:fullperiod}) is uniform in $f \ge 2$. The First-G
localization $k^{\star}(b) = \spf(b)$ (Theorem~\ref{thm:firstg})
identifies the smallest non-unit residue index. The spectral
characterization (Theorem~\ref{thm:obstruction-zero}) and the
synchronization at the smallest spectral zero
(Theorem~\ref{thm:sync}) connect the obstruction sequence to the
zero set of the product $f_{b}(k) = \prod_{p \mid b} R(k, p)$. The
layered-divisor structure (Theorem~\ref{thm:layered}) on squarefree
$b$ gives the exact count $2^{j} - 1$ of non-trivial divisors of $b$
whose Fej\'{e}r quotient vanishes at the $j$-th primorial $b_{j}$.
The continuum limit (Theorem~\ref{thm:limit}) reads $R(k, f)$ as a
discrete approximant to $\sinc^{2}(k/f)$, and the asymptotic average
across the corridor $\{1, \dots, f-1\}$ converges to the sine-integral
constant $\Si(2\pi)/\pi \approx 0.4514$
(Theorem~\ref{thm:asymptotic}).

\paragraph{Positioning and contribution.} The closed form for $R$
(Theorem~\ref{thm:closedform}) and the continuum limit
(Theorem~\ref{thm:limit}) are classical Fej\'{e}r-kernel material,
included to make the paper self-contained. The full-period
cancellation (Theorem~\ref{thm:fullperiod}) and the First-G
localization (Theorem~\ref{thm:firstg}) follow each from one line of
elementary algebra. The substantive contributions are
\emph{(i)} the spectral characterization of the obstruction sequence
(Theorem~\ref{thm:obstruction-zero}): for every $b > 1$ with distinct
prime factors $p_{1}, \dots, p_{r}$, the spectral product
$f_{b}(k) = \prod_{j} R(k, p_{j})$ acts as a continuous-in-$k$
indicator function for $\gcd(k, b) > 1$, with the divisibility lattice
of $\rad(b)$ entering directly through the zero loci of the factors;
\emph{(ii)} the layered-divisor count $2^{j} - 1$ at the $j$-th
primorial divisor (Theorem~\ref{thm:layered}), built on the Boolean
divisor lattice of squarefree $b$ and the uniform full-period
cancellation; and \emph{(iii)} the closed-form evaluation of the
corridor average as $\Si(2\pi)/\pi$ via Riemann-sum
$\to \int_{0}^{1} \sinc^{2}(t)\, dt$ followed by integration by parts
(Theorem~\ref{thm:asymptotic}). The asymptotic zero density of $f_{b}$
(Corollary~\ref{cor:asymp-density}) recovers the Euler product
$1 - \varphi(\rad(b))/\rad(b)$ as a corollary of the spectral
characterization. Related classical work on the distribution of the
smallest prime factor~\cite{Erdos1959,Pomerance1985,Tenenbaum2015}
and sieve-theoretic treatments~\cite{IwaniecKowalski2004,FriedlanderIwaniec2010}
provide the analytic-number-theory context.

\paragraph{Tier discipline.} The seven theorems all PROVE the claims
asserted; we COMPUTE verifications via two scripts at exact integer
arithmetic or machine-precision floating point (Section~\ref{sec:verify}).
The exact identity $\sinc^{2}(1/2) = (2/3)/\zeta(2)$ at the corridor
midpoint is one-line from $\zeta(2) = \pi^{2}/6$ and is cited as
STRUCTURAL RHYME, not theorem. The corridor-midpoint substrate
question (why $t = 1/2$ makes this identity structurally relevant) is
OPEN, as is the question of whether $f_{b}$ admits a Ramanujan-sum or
character-sum form beyond the inclusion-exclusion identity in
Corollary~\ref{cor:incl-excl}.

\paragraph{Organization.} Section~\ref{sec:setup} gives the setup.
Section~\ref{sec:fullperiod} establishes the closed form
(Theorem~\ref{thm:closedform}) and full-period cancellation
(Theorem~\ref{thm:fullperiod}). Section~\ref{sec:firstg} proves the
First-G localization (Theorem~\ref{thm:firstg}).
Section~\ref{sec:spectral} establishes the spectral characterization
(Theorem~\ref{thm:obstruction-zero}), records the synchronization
(Theorem~\ref{thm:sync}) as a special case, and derives the
inclusion-exclusion identity and asymptotic zero density
(Corollaries~\ref{cor:incl-excl},~\ref{cor:asymp-density}).
Section~\ref{sec:layered} proves the layered-divisor structure
(Theorem~\ref{thm:layered}). Section~\ref{sec:asymptotic} proves the
continuum limit (Theorem~\ref{thm:limit}) and the corridor-average
identity (Theorem~\ref{thm:asymptotic}). Section~\ref{sec:verify}
reports the verification. Section~\ref{sec:scope} lists scope and
limitations.

%% -------- section 2: setup --------
\section{Setup and notation}\label{sec:setup}

Throughout, $b > 1$ is an integer with prime factorization
$b = p_{1}^{a_{1}} \cdots p_{r}^{a_{r}}$, distinct primes
$p_{1} < p_{2} < \cdots < p_{r}$, $\spf(b) = p_{1}$, and
$\rad(b) = p_{1} p_{2} \cdots p_{r}$. The integer $b$ is
\emph{squarefree} iff $b = \rad(b)$, equivalently every $a_{i} = 1$.
The normalized sinc function is
\[
  \sinc(x) =
  \begin{cases}
    \sin(\pi x)/(\pi x) & x \ne 0 \\
    1 & x = 0
  \end{cases},
  \qquad
  \sinc^{2}(x) = \sinc(x)^{2}.
\]
The discrete Fej\'{e}r quotient $R(k, f)$ is defined in~\eqref{eq:Rdef}.

\begin{definition}[Coprimality partition; First-G event]\label{def:partition}
For $b > 1$ and $k \ge 1$, the \emph{coprimality partition} of the
alphabet $\{1, \dots, k\}$ relative to $b$ is the ordered pair
\[
  C_{k}(b) = \{x \in \{1, \dots, k\} : \gcd(x, b) = 1\},
  \qquad
  G_{k}(b) = \{1, \dots, k\} \setminus C_{k}(b).
\]
The \emph{First-G event} of $b$ is the integer
$k^{\star}(b) := \min\{k \ge 1 : |G_{k}(b)| > 0\}$.
\end{definition}

\begin{definition}[Spectral product]\label{def:spectral-product}
For $b > 1$ with distinct prime factors $p_{1} < \cdots < p_{r}$, the
\emph{spectral product} of $b$ is
\begin{equation}\label{eq:spectral-product}
  f_{b}(k) \;:=\; \prod_{j=1}^{r} R(k, p_{j}), \qquad k \ge 1.
\end{equation}
\end{definition}

By~\eqref{eq:Rdef}, $f_{b}$ has the closed form
\[
  f_{b}(k) \;=\; \frac{1}{k^{2r}}
                 \prod_{j=1}^{r}
                 \frac{\sin^{2}(\pi k/p_{j})}{\sin^{2}(\pi/p_{j})}.
\]

%% -------- section 3: full-period cancellation --------
\section{The closed form and full-period cancellation}\label{sec:fullperiod}

\begin{theorem}[Discrete Fej\'{e}r identity]\label{thm:closedform}
For every integer $f \ge 2$ and every integer $k \ge 1$,
\begin{equation}\label{eq:closedform}
  R(k, f) \;=\; \frac{\sin^{2}(\pi k/f)}{k^{2}\,\sin^{2}(\pi/f)}.
\end{equation}
\end{theorem}

\begin{proof}
The geometric-series formula gives
\[
  \sum_{j=1}^{k} e^{2\pi i j/f}
  \;=\;
  e^{2\pi i/f} \cdot \frac{1 - e^{2\pi i k/f}}{1 - e^{2\pi i/f}}.
\]
Taking the squared modulus of $(1/k)$ times this sum and applying
$|1 - e^{i\theta}|^{2} = 4 \sin^{2}(\theta/2)$,
\[
  R(k, f)
  = \frac{1}{k^{2}}\,
    \frac{|1 - e^{2\pi i k/f}|^{2}}{|1 - e^{2\pi i/f}|^{2}}
  = \frac{1}{k^{2}}\,
    \frac{4 \sin^{2}(\pi k/f)}{4 \sin^{2}(\pi/f)}
  = \frac{\sin^{2}(\pi k/f)}{k^{2} \sin^{2}(\pi/f)}.
\qedhere
\]
\end{proof}

\begin{theorem}[Full-period cancellation]\label{thm:fullperiod}
For every $f \ge 2$ and every $k \ge 1$,
\[
  R(k, f) = 0 \;\Longleftrightarrow\; f \mid k.
\]
\end{theorem}

\begin{proof}
By~\eqref{eq:closedform}, the denominator $k^{2} \sin^{2}(\pi/f)$ is
strictly positive for every $k \ge 1$ and $f \ge 2$, so the zero set
of $R(\cdot, f)$ is determined entirely by the numerator
$\sin^{2}(\pi k/f)$. Now $\sin^{2}(\pi k/f) = 0 \Leftrightarrow
\sin(\pi k/f) = 0 \Leftrightarrow \pi k/f \in \pi \Z \Leftrightarrow
k/f \in \Z \Leftrightarrow f \mid k$.
\end{proof}

Theorem~\ref{thm:fullperiod} is the \emph{full-period} statement: the
zeros of $R(\cdot, f)$ form the arithmetic progression
$\{f, 2f, 3f, \dots\}$ for every $f \ge 2$, prime or composite.

\begin{corollary}[Endpoint values and monotonicity]\label{cor:endpoints}
For every $f \ge 2$:
\begin{enumerate}[label={\upshape(\roman*)}]
\item $R(1, f) = 1$ (full resonance at $k = 1$);
\item $R(f - 1, f) = 1/(f - 1)^{2}$, the last positive value before
      the first zero;
\item $R(f, f) = 0$ (exact zero at $k = f$, per
      Theorem~\ref{thm:fullperiod});
\item $R(k, f)$ is strictly decreasing on $\{1, \dots, f - 1\}$.
\end{enumerate}
\end{corollary}

\begin{proof}
\emph{(i)} Direct from~\eqref{eq:closedform}: at $k = 1$,
$R(1, f) = \sin^{2}(\pi/f) / \sin^{2}(\pi/f) = 1$.
\emph{(iii)} At $k = f$, $\sin(\pi) = 0$.
\emph{(ii)} At $k = f - 1$, $\sin(\pi(f-1)/f) = \sin(\pi - \pi/f)
= \sin(\pi/f)$, so $R(f-1, f) = \sin^{2}(\pi/f) / ((f-1)^{2}
\sin^{2}(\pi/f)) = 1/(f-1)^{2}$.
\emph{(iv)} Factor~\eqref{eq:closedform} as
$R(k, f) = (\pi/f)^{2} \sin^{-2}(\pi/f) \cdot (\sin(\pi k/f)/(\pi k/f))^{2}$.
The first factor is constant in $k$. For $1 \le k \le f - 1$,
$\pi k/f \in (0, \pi)$, and the function $x \mapsto \sin(x)/x$ is
strictly decreasing on $(0, \pi)$ (its derivative
$(x \cos x - \sin x)/x^{2}$ is negative on $(0, \pi)$ via
$h(x) := \sin x - x \cos x$, $h(0) = 0$, $h'(x) = x \sin x > 0$ on
$(0, \pi)$). The square of a positive strictly decreasing function
is strictly decreasing.
\end{proof}

%% -------- section 4: First-G localization --------
\section{First-G event localization}\label{sec:firstg}

The number-theoretic side of the synchronization is the First-G
localization theorem: the First-G event of $b$ is the smallest prime
factor of $b$. The proof is one line of $\gcd$ reasoning and does not
require squarefree-ness.

\begin{theorem}[First-G localization]\label{thm:firstg}
Let $b > 1$ be an integer with smallest prime factor $p_{1} = \spf(b)$.
Then $k^{\star}(b) = p_{1}$. More precisely:
\begin{enumerate}[label={\upshape(\roman*)}]
\item $|G_{k}(b)| = 0$ for every $k$ with $1 \le k < p_{1}$, and
\item $|G_{p_{1}}(b)| = 1$, with $G_{p_{1}}(b) = \{p_{1}\}$.
\end{enumerate}
\end{theorem}

\begin{proof}
\emph{(i)} Fix $k < p_{1}$ and $x \in \{1, \dots, k\}$. If
$\gcd(x, b) = d > 1$, then $d$ has a prime divisor $q$ with $q \mid b$;
hence $q \ge p_{1}$ by definition of $p_{1}$. But $q \mid x$ forces
$q \le x \le k < p_{1}$, a contradiction. Therefore $\gcd(x, b) = 1$
and $G_{k}(b) = \emptyset$.
\emph{(ii)} At $k = p_{1}$, the element $p_{1}$ enters the alphabet,
and $\gcd(p_{1}, b) = p_{1} > 1$ since $p_{1} \mid b$, so
$p_{1} \in G_{p_{1}}(b)$. For any $x \in \{1, \dots, p_{1}\}$ with
$x \ne p_{1}$ we have $x < p_{1}$, so part~(i) gives $\gcd(x, b) = 1$.
\end{proof}

\begin{remark}\label{rem:nonsqfree}
Theorem~\ref{thm:firstg} uses only the definitions of $p_{1}$ and
$\gcd$; squarefree-ness of $b$ is not required. For example, at
$b = 12$, $\spf(b) = 2$ and $G_{2}(12) = \{2\}$ in the same way as
$G_{2}(6) = \{2\}$. The verification sample below restricts to
squarefree $b$ for convenience only, as that is the regime where the
companion layered-closure theorem (Theorem~\ref{thm:layered}) earns
its keep.
\end{remark}

%% -------- section 5: spectral characterization --------
\section{Spectral characterization of the obstruction sequence}\label{sec:spectral}

The synchronization at $k = p_{1}$ extends to the entire obstruction
sequence via the spectral product $f_{b}$ of
Definition~\ref{def:spectral-product}.

\begin{theorem}[Obstruction--zero correspondence]\label{thm:obstruction-zero}
Let $b > 1$ with distinct prime factors $p_{1} < \cdots < p_{r}$,
and let $f_{b}$ be the spectral product of
Definition~\ref{def:spectral-product}. Then for every integer
$k \ge 1$,
\begin{equation}\label{eq:obstruction-zero}
  f_{b}(k) = 0
  \;\Longleftrightarrow\;
  \gcd(k, b) > 1
  \;\Longleftrightarrow\;
  k \in G_{k}(b).
\end{equation}
Equivalently, the zero set of $f_{b}$ in $\N$ is
\[
  Z(f_{b}) \;=\; \bigcup_{j=1}^{r} p_{j}\,\N,
\]
the union of the arithmetic progressions of multiples of each
distinct prime factor of $b$.
\end{theorem}

\begin{proof}
By Theorem~\ref{thm:fullperiod}, $R(k, p_{j}) = 0 \Leftrightarrow
p_{j} \mid k$. The product $f_{b}(k)$ vanishes iff some factor
$R(k, p_{j})$ vanishes iff $p_{j} \mid k$ for some $j$ iff
$\gcd(k, b) > 1$. The final equivalence is Definition~\ref{def:partition}.
\end{proof}

\begin{theorem}[Synchronization at the smallest spectral zero]\label{thm:sync}
For every integer $b > 1$ with $p_{1} = \spf(b)$, the First-G event
$k^{\star}(b)$ and the smallest integer zero of $f_{b}$ coincide at
$k = p_{1}$:
\[
  k^{\star}(b)
  \;=\;
  \min\{k \ge 1 : f_{b}(k) = 0\}
  \;=\;
  p_{1}.
\]
\end{theorem}

\begin{proof}
By Theorem~\ref{thm:firstg}, $k^{\star}(b) = p_{1}$. By
Theorem~\ref{thm:obstruction-zero}, the zero set of $f_{b}$ is
$\bigcup_{j} p_{j} \N$; its smallest element is $p_{1}$, the smallest
prime factor of $b$. Both quantities equal $p_{1}$.
\end{proof}

\begin{remark}[No squarefree restriction; $f_{b}$ depends only on $\rad(b)$]\label{rem:radical}
Theorem~\ref{thm:obstruction-zero} holds for every integer $b > 1$,
not only for squarefree $b$. Definition~\ref{def:spectral-product}
forms $f_{b}(k)$ from the \emph{distinct} prime factors of $b$, which
are precisely the prime factors of $\rad(b)$, so
$f_{b}(k) = f_{\rad(b)}(k)$ for every $k$. For example,
$b = 12 = 2^{2} \cdot 3$ has distinct primes $\{2, 3\}$, so
$f_{12}(k) = R(k, 2)\, R(k, 3) = f_{6}(k)$, and the zero set
$\bigcup\{2\N, 3\N\}$ is the set of $k$ with $\gcd(k, 12) > 1$.
Corollary~\ref{cor:asymp-density} below is similarly stated for
general $b > 1$.
\end{remark}

\begin{corollary}[Inclusion--exclusion identity for $|G_{k}|$]\label{cor:incl-excl}
With $b$ and $f_{b}$ as in Theorem~\ref{thm:obstruction-zero}, the
number of obstructions in $\{1, \dots, k\}$ is
\[
  |G_{k}(b)|
  \;=\; \#\{j \in \{1, \dots, k\} : f_{b}(j) = 0\}
  \;=\; -\sum_{\substack{d \mid \rad(b)\\ d > 1}}
              \mu(d)\,\lfloor k/d \rfloor.
\]
\end{corollary}

\begin{proof}
The first equality is Theorem~\ref{thm:obstruction-zero}. The second
is the standard inclusion-exclusion count of integers in $[1, k]$
divisible by at least one $p_{j}$, applied to the squarefree integer
$\rad(b)$ (whose divisors $d$ are squarefree products of the $p_{j}$,
$\mu(d) = (-1)^{\omega(d)}$).
\end{proof}

\begin{corollary}[Asymptotic zero density of $f_{b}$]\label{cor:asymp-density}
For every $b > 1$ with distinct prime factors $p_{1}, \dots, p_{r}$,
\begin{equation}\label{eq:asymp-density}
  \lim_{K \to \infty}
  \frac{1}{K}\,\#\{j \in \{1, \dots, K\} : f_{b}(j) = 0\}
  \;=\; 1 - \prod_{j=1}^{r}\!\left(1 - \frac{1}{p_{j}}\right)
  \;=\; 1 - \frac{\varphi(\rad(b))}{\rad(b)}.
\end{equation}
\end{corollary}

\begin{proof}
By Theorem~\ref{thm:obstruction-zero} and Corollary~\ref{cor:incl-excl},
the LHS equals $\lim_{K \to \infty} |G_{K}(b)|/K$, which by the Euler
product for the multiplicative density of integers coprime to
$\rad(b)$~\cite{HardyWright2008,Apostol1976} equals
$1 - \prod_{j}(1 - 1/p_{j})$.
\end{proof}

%% -------- section 6: layered-divisor structure --------
\section{Layered-divisor structure on squarefree moduli}\label{sec:layered}

Theorem~\ref{thm:obstruction-zero} states that $f_{b}(k) = 0$ iff
\emph{some} prime factor of $b$ divides $k$. The layered-divisor
structure of this section refines the observation to a count: at the
$j$-th primorial divisor of squarefree $b$, exactly $2^{j} - 1$ of the
non-trivial divisors of $b$ produce $R(k, d) = 0$.

For squarefree $b = p_{1} p_{2} \cdots p_{r}$ with
$p_{1} < \cdots < p_{r}$, the divisor lattice is the Boolean lattice
on $\{p_{1}, \dots, p_{r}\}$: divisors of $b$ are exactly the products
$\prod_{i \in S} p_{i}$ for subsets $S \subseteq \{1, \dots, r\}$, with
the non-trivial divisors ($> 1$) corresponding to non-empty subsets.

\begin{theorem}[Squarefree layered-divisor structure]\label{thm:layered}
Let $b = p_{1} p_{2} \cdots p_{r}$ be squarefree with
$p_{1} < \cdots < p_{r}$. The set
\[
  Z(b)
  \;:=\;
  \{k \in \{1, \dots, b\} : R(k, d) = 0
   \text{ for at least one non-trivial divisor } d \mid b\}
\]
satisfies the following three statements.
\begin{enumerate}[label={\upshape(\roman*)}]
\item $Z(b) = \{k \in \{1, \dots, b\} : \gcd(k, b) > 1\}$.
\item The smallest element of $Z(b)$ is $k = p_{1} = \spf(b)$.
\item For each $j \in \{1, \dots, r\}$, set $b_{j} = p_{1} p_{2}
\cdots p_{j}$ (the $j$-th primorial divisor of $b$). The number of
non-trivial divisors $d \mid b$ for which $R(b_{j}, d) = 0$ is exactly
$2^{j} - 1$.
\end{enumerate}
\end{theorem}

\begin{proof}
\emph{(i)} By Theorem~\ref{thm:fullperiod}, $R(k, d) = 0 \Leftrightarrow
d \mid k$. So $k \in Z(b)$ iff some non-trivial divisor $d \mid b$
satisfies $d \mid k$, iff $\gcd(k, b)$ has a non-trivial divisor, iff
$\gcd(k, b) > 1$.
\emph{(ii)} By~(i), the smallest element of $Z(b)$ is the smallest
$k \in \{1, \dots, b\}$ with $\gcd(k, b) > 1$, which by
Theorem~\ref{thm:firstg} is $\spf(b) = p_{1}$.
\emph{(iii)} Fix $j$ and write $b_{j} = p_{1} \cdots p_{j}$. A
non-trivial divisor $d \mid b$ is of the form $d = \prod_{i \in S} p_{i}$
for some non-empty subset $S \subseteq \{1, \dots, r\}$. By
Theorem~\ref{thm:fullperiod}, $R(b_{j}, d) = 0 \Leftrightarrow d \mid b_{j}$.
A divisor $d = \prod_{i \in S} p_{i}$ divides $b_{j} = \prod_{i=1}^{j}
p_{i}$ iff $S \subseteq \{1, \dots, j\}$ (here we use squarefree-ness
of $b$, which makes the divisor lattice Boolean). The number of
non-empty subsets of $\{1, \dots, j\}$ is $2^{j} - 1$.
\end{proof}

\begin{remark}[Non-squarefree extension]\label{rem:layered-nonsf}
If $b$ is not squarefree, write $b = \rad(b) \cdot m$ with $m \ge 1$.
Theorem~\ref{thm:layered} then applies to $\rad(b)$, and the count
$2^{j} - 1$ at the $j$-th primorial divisor of $\rad(b)$ remains
exact. The divisor lattice of $b$ itself acquires multiplicities, and
the full non-squarefree count appears not to admit a closed form as
compact as $2^{j} - 1$.
\end{remark}

%% -------- section 7: continuum + asymptotic --------
\section{Continuum and asymptotic limits}\label{sec:asymptotic}

The harmonic side of the picture is completed by two limit statements:
a pointwise continuum limit identifying $R(k, f)$ with a discrete
approximant of $\sinc^{2}(k/f)$, and an averaged limit identifying the
corridor mean of $R(\cdot, f)$ with $\Si(2\pi)/\pi$.

\begin{theorem}[Continuum limit; rectangular-pulse spectrum]\label{thm:limit}
For every $t \in \R$ with $t > 0$,
\[
  \lim_{f \to \infty,\ k/f \to t} R(k, f) \;=\; \sinc^{2}(t).
\]
\end{theorem}

\begin{proof}
From~\eqref{eq:closedform}, $R(k, f) = \sin^{2}(\pi k/f) /
(k^{2} \sin^{2}(\pi/f))$. As $f \to \infty$ with $k/f \to t$ fixed,
$k \cdot \sin(\pi/f) = (\pi/f) \cdot k \cdot (1 + O(1/f^{2}))
\to \pi t$, so $k^{2} \sin^{2}(\pi/f) \to \pi^{2} t^{2}$. The
numerator satisfies $\sin^{2}(\pi k/f) \to \sin^{2}(\pi t)$. Hence
$R(k, f) \to \sin^{2}(\pi t)/(\pi^{2} t^{2}) = \sinc^{2}(t)$.
\end{proof}

\begin{remark}[Specific exact values; structural rhyme]\label{rem:exactvalues}
Two exact evaluations of $\sinc^{2}$ that arise as continuum limits of
$R(k, f)$ at rational arguments:
\begin{align*}
  \sinc^{2}(1/2) &= \frac{4}{\pi^{2}} = \frac{2/3}{\zeta(2)}, \\
  \sinc^{2}(1/10) &= \frac{25(\sqrt{5}-1)^{2}}{4\pi^{2}}.
\end{align*}
The first follows from $\sin(\pi/2) = 1$ together with
$\zeta(2) = \pi^{2}/6$; the second from $\sin(\pi/10) = (\sqrt{5}-1)/4$.
Both are cited as structural rhymes, not theorems.
\end{remark}

\begin{theorem}[Asymptotic corridor average]\label{thm:asymptotic}
\[
  \lim_{f \to \infty}
  \frac{1}{f - 1} \sum_{k=1}^{f-1} R(k, f)
  \;=\;
  \int_{0}^{1} \sinc^{2}(t)\, dt
  \;=\;
  \frac{\Si(2\pi)}{\pi}
  \;\approx\; 0.4514,
\]
where $\Si(x) = \int_{0}^{x} \sin(t)/t\, dt$.
\end{theorem}

\begin{proof}
By Theorem~\ref{thm:limit}, $R(k, f) \to \sinc^{2}(k/f)$ pointwise as
$f \to \infty$ with $k/f$ fixed. The sum
\[
  \frac{1}{f - 1} \sum_{k=1}^{f-1} R(k, f)
  \;=\;
  \frac{1}{f - 1} \sum_{k=1}^{f-1} \sinc^{2}(k/f)
  + O\!\left(\frac{1}{f^{2}}\right)
\]
is a Riemann sum for $\int_{0}^{1} \sinc^{2}(t)\, dt$ on the partition
$\{1/f, 2/f, \dots, (f-1)/f\}$ of step $1/f$. Since $\sinc^{2}$ is
$C^{\infty}$ on $[0, 1]$ and the $O(1/f^{2})$ pointwise correction
vanishes uniformly in $k$, the Riemann sum converges to the integral.

The integral evaluates by integration by parts. Substituting
$u = \pi t$,
\[
  \int_{0}^{1} \sinc^{2}(t)\, dt
  = \frac{1}{\pi} \int_{0}^{\pi} \frac{\sin^{2}(u)}{u^{2}}\, du.
\]
With $dv = du/u^{2}$ (so $v = -1/u$), $w = \sin^{2}(u)$,
$dw = \sin(2u)\, du$,
\[
  \int_{0}^{\pi} \frac{\sin^{2}(u)}{u^{2}}\, du
  = \left[ -\frac{\sin^{2}(u)}{u} \right]_{0}^{\pi}
  + \int_{0}^{\pi} \frac{\sin(2u)}{u}\, du
  = 0 + \int_{0}^{2\pi} \frac{\sin(s)}{s}\, ds
  = \Si(2\pi),
\]
where the boundary term vanishes at $u = 0$ via $\sin^{2}(u)/u = O(u)$
and at $u = \pi$ via $\sin(\pi) = 0$, and we substituted $s = 2u$ in
the remaining integral (using $s$ rather than $t$ to avoid reusing
the outer integration variable from~$\int_{0}^{1}\sinc^{2}(t)\, dt$).
Hence $\int_{0}^{1} \sinc^{2}(t)\, dt = \Si(2\pi)/\pi$.
\end{proof}

\begin{remark}[Numerical value]\label{rem:siNumerical}
$\Si(2\pi) = 1.41815\dots$, so $\Si(2\pi)/\pi = 0.45141\dots$. The
corridor-midpoint value $\sinc^{2}(1/2) = 4/\pi^{2} = 0.40528\dots$
sits strictly below the corridor average, as expected since
$\sinc^{2}$ is concave on a neighborhood of $1/2$.
\end{remark}

%% -------- section 8: verification --------
\section{Verification}\label{sec:verify}

All claims are verified by two scripts archived at Zenodo bundle
DOI~\texttt{10.5281/zenodo.18852047}: \texttt{proof\_first\_g\_event.py}
(exhaustive integer enumeration for Theorem~\ref{thm:firstg}) and
\texttt{verify\_J03.py} (machine-precision checks for the remaining
theorems and corollaries). Total runtime is under five seconds on a
2024 consumer laptop.

\begin{enumerate}[label={\upshape(\arabic*)}]
\item \textbf{Theorem~\ref{thm:firstg} (First-G localization).} For
every squarefree $b$ in $[2, 500]$ (305 values), $k^{\star}(b)$ is
computed by direct $\gcd$ enumeration and matched against $\spf(b)$.
Total: 22{,}367 $(b, k)$ pairs, zero counterexamples.

\item \textbf{Theorem~\ref{thm:closedform} (closed form).} For every
prime $f \in \{3, 5, 7, 11, 13, 17, 19, 23, 29, 31, 37, 41, 43, 47\}$
and every $k \in \{1, \dots, f + 1\}$, the closed form is checked
against the literal geometric sum~\eqref{eq:Rgeom}. Maximum deviation
$4.44 \times 10^{-16}$, i.e., machine epsilon.

\item \textbf{Theorem~\ref{thm:fullperiod} (full-period cancellation).}
For every prime $p \in \{3, 5, 7, 11, 13, \dots, 199\}$ (45 primes) and
every $k \in \{1, \dots, p\}$, $R(k, p) = 0 \Leftrightarrow p \mid k$
is verified via exact integer divisibility cross-checked with the
floating-point evaluation. Zero counterexamples in 4{,}225 $(p, k)$
pairs.

\item \textbf{Theorem~\ref{thm:sync} (synchronization).} For
$b \in \{10, 30, 35, 77, 105, 210, 1001, 2310\}$, the equality
$k^{\star}(b) = \spf(b)$ is matched against $R(\spf(b), \spf(b)) = 0$
to machine precision.

\item \textbf{Theorem~\ref{thm:obstruction-zero} (obstruction--zero
correspondence).} For every squarefree $b \in [2, 50]$ (30 values)
and every $k \in [1, 30]$, the boolean equivalence
$f_{b}(k) = 0 \Leftrightarrow \gcd(k, b) > 1$ is verified at exact
integer arithmetic via the factor-level test
$|\sin(\pi k/p_{j})| < 10^{-12}$. Total: 900 cells, 900/900 matches.

\item \textbf{Corollary~\ref{cor:asymp-density} (asymptotic density).}
For $b \in \{2, 6, 10, 15, 30, 42, 105, 210, 2310\}$, the integer
obstruction count $(1/K)\,|G_{K}(b)|$ at $K = 100{,}000$ is matched
against the Euler product $1 - \varphi(\rad(b))/\rad(b)$. Maximum
deviation: $6.67 \times 10^{-6}$.

\item \textbf{Theorem~\ref{thm:layered} (squarefree layered closure).}
For 50 squarefree $b \in \{6, 10, 14, 15, 21, 22, \dots, 210\}$: the
smallest $k$ at which any non-trivial divisor $d \mid b$ satisfies
$R(k, d) = 0$ is exactly $\spf(b)$; the count at the second primorial
$b_{2} = p_{1} p_{2}$ is exactly $2^{2} - 1 = 3$; for $b$ with
$\omega(b) \ge 3$, the count at the third primorial $b_{3}$ is
exactly $2^{3} - 1 = 7$. Zero counterexamples.

\item \textbf{Theorem~\ref{thm:limit} (continuum limit).} Scaling
$f \in \{10, 100, 1000, 10000\}$ at $t = 1/2$ and $f \in \{4, 40, 400,
4000\}$ at $t = 1/4$, the difference $|R(k, f) - \sinc^{2}(t)|$ decays
as $O(1/f^{2})$ and is below $2 \times 10^{-9}$ at $f = 10{,}000$.

\item \textbf{Theorem~\ref{thm:asymptotic} (corridor average).}
Numerical Riemann sums at $f \in \{50, 100, 500, 1000\}$ converge to
$\Si(2\pi)/\pi \approx 0.45141$:
\begin{center}
\begin{tabular}{c|c|c|c}
\toprule
$f$ & $\frac{1}{f-1}\sum_{k=1}^{f-1} R(k, f)$ & $\Si(2\pi)/\pi$ & deviation \\
\midrule
$50$   & $0.4510$ & $0.4514$ & $4.0 \times 10^{-4}$ \\
$100$  & $0.4511$ & $0.4514$ & $3.4 \times 10^{-4}$ \\
$500$  & $0.4513$ & $0.4514$ & $9.2 \times 10^{-5}$ \\
$1000$ & $0.4514$ & $0.4514$ & $4.8 \times 10^{-5}$ \\
\bottomrule
\end{tabular}
\end{center}
The convergence is consistent with the Riemann-sum error term and the
$O(1/f^{2})$ pointwise correction; deviation drops by an order of
magnitude as $f$ ranges over $\{50, 100, 500, 1000\}$.
\end{enumerate}

%% -------- section 9: scope --------
\section{Scope and limitations}\label{sec:scope}

\begin{enumerate}[label={\upshape(\arabic*)}]
\item \textbf{What is proved.} Seven theorems
(\ref{thm:closedform}--\ref{thm:asymptotic}) and two corollaries
(\ref{cor:incl-excl},~\ref{cor:asymp-density}): closed form, full-period
cancellation, First-G localization, obstruction-zero correspondence,
synchronization, layered-divisor structure, continuum limit, and
corridor average. The proofs are elementary throughout.

\item \textbf{What is verified.} Two scripts, nine checks, exact
arithmetic where possible and machine-precision floating point
otherwise. Zero counterexamples in the integer-arithmetic portions;
deviation below $5 \times 10^{-5}$ at $f = 1000$ for the asymptotic
average; deviation $6.67 \times 10^{-6}$ at $K = 100{,}000$ for the
density corollary.

\item \textbf{Squarefree restriction.} Closed form, full-period
cancellation, First-G localization, obstruction-zero correspondence,
and synchronization are uniform in $b$ and do not require
squarefree-ness; for $f_{b}$, only the distinct prime factors of $b$
enter (Remark~\ref{rem:radical}). The squarefree restriction
genuinely enters only in Theorem~\ref{thm:layered}, where the divisor
lattice is Boolean and the count $2^{j} - 1$ is exact. For
non-squarefree $b$, the layered count at the $j$-th primorial divisor
of $\rad(b)$ remains $2^{j} - 1$, but the count over divisors of $b$
itself acquires multiplicities (Remark~\ref{rem:layered-nonsf}).

\item \textbf{What is \emph{not} claimed.} No statement about the
analytic continuation of $\zeta(s)$; no consequence for the Riemann
hypothesis; no implication for primes in arithmetic progressions; no
polynomial-time factoring algorithm. The identity $\sinc^{2}(1/2) =
(2/3)/\zeta(2)$ is one line from $\zeta(2) = \pi^{2}/6$ and is cited
as structural rhyme only.

\item \textbf{Open questions.} (a)~Why does the corridor midpoint at
$t = 1/2$ make $\sinc^{2}(1/2) = (2/3)/\zeta(2)$ structurally
relevant beyond its one-line derivation? (b)~Does the spectral product
$f_{b}$ admit a Ramanujan-sum or character-sum form beyond the
inclusion-exclusion identity of Corollary~\ref{cor:incl-excl}? (c)~Is
there a closed-form non-squarefree count refining
Remark~\ref{rem:layered-nonsf}?

\item \textbf{Finite-range verification.} The exhaustive integer
checks are on the stated finite ranges. The theorems themselves are
proved in full generality; the finite verification is reassurance,
not a pillar of the result.
\end{enumerate}

%% -------- references --------
\begin{thebibliography}{99}

\bibitem{Apostol1976}
T.~M.~Apostol.
\newblock \emph{Introduction to Analytic Number Theory}.
\newblock Undergraduate Texts in Mathematics. Springer, New York, 1976.

\bibitem{Erdos1959}
P.~Erd\H{o}s.
\newblock On the distribution of the smallest prime factor of $n$.
\newblock \emph{Mathematika}, 6:1--6, 1959.

\bibitem{Fejer1900}
L.~Fej\'{e}r.
\newblock Sur les fonctions born\'{e}es et int\'{e}grables.
\newblock \emph{Comptes Rendus de l'Acad\'{e}mie des Sciences, Paris},
131:984--987, 1900.

\bibitem{FriedlanderIwaniec2010}
J.~Friedlander and H.~Iwaniec.
\newblock \emph{Opera de Cribro}.
\newblock American Mathematical Society Colloquium Publications 57.
\newblock American Mathematical Society, Providence, RI, 2010.

\bibitem{HardyWright2008}
G.~H.~Hardy and E.~M.~Wright.
\newblock \emph{An Introduction to the Theory of Numbers}, 6th edition.
\newblock Revised by D.~R.~Heath-Brown and J.~H.~Silverman.
\newblock Oxford University Press, Oxford, 2008.

\bibitem{IwaniecKowalski2004}
H.~Iwaniec and E.~Kowalski.
\newblock \emph{Analytic Number Theory}.
\newblock American Mathematical Society Colloquium Publications 53.
\newblock American Mathematical Society, Providence, RI, 2004.

\bibitem{Pomerance1985}
C.~Pomerance.
\newblock The role of smooth numbers in number-theoretic algorithms.
\newblock \emph{Proc.\ Int.\ Congr.\ Math.}\ (Z\"{u}rich, 1994),
volume~1, pages 411--422. Birkh\"{a}user, 1995.

\bibitem{Tenenbaum2015}
G.~Tenenbaum.
\newblock \emph{Introduction to Analytic and Probabilistic Number Theory},
3rd edition.
\newblock Cambridge Studies in Advanced Mathematics 46. Cambridge
University Press, Cambridge, 2015.

\bibitem{Zygmund2002}
A.~Zygmund.
\newblock \emph{Trigonometric Series}, 3rd edition.
\newblock Cambridge University Press, Cambridge, 2002.

%% -- companion papers ------------------------
\bibitem{SandersGish2026Sigma}
B.~R.~Sanders and M.~Gish.
\newblock Non-Associativity Decay in Binary Composition Tables over
$\mathbb{Z}/N\mathbb{Z}$.
\newblock Preprint, 2026; J-series J14.
\newblock Archived in shared deposit DOI 10.5281/zenodo.18852047.

\bibitem{SandersGish2026FourCore}
B.~R.~Sanders and M.~Gish.
\newblock Joint Closure, Per-Coordinate Fuse Data, and a Closed-Form
Algebraic Attractor of Two Commutative Binary Operations on
$\mathbb{Z}/10\mathbb{Z}$.
\newblock Preprint, 2026; J-series J15.
\newblock Archived in shared deposit DOI 10.5281/zenodo.18852047.

\end{thebibliography}

\end{document}
